\section{Repaso y reflexión de la sección}

\paragraph{}
Los primeros 15 minutos de la Lecture 05 delimitaron el dominio de planning domain mediante curvas de costo computacional y desempeño: la ``smooth curve'' del LLM ya mostraba rendimientos marginales decrecientes del scaling, por lo que era necesario buscar una ruta de solución más estructurada.

\paragraph{}
Search Former y DU Forer convierten el search trace en una secuencia con supervision, y usan dropout para simular la información de búsqueda incompleta que existe en la realidad. Los benchmarks de Soang y Maze mostraron que el modelo de trace puede, con un modelo pequeño, superar al baseline solution-only en generalization, y ofrecer recomendaciones más consistentes para el multi-turn agent.

\paragraph{}
La parte de interacción de múltiples turnos forma, mediante la APAC Constitution y el DPO loop del agent judge, un ciclo autosupervisado que enfatiza la medición en cuatro dimensiones y una estrategia de preguntas diferenciadas por persona, lo que permite que el agent produzca un plan estructurado dentro de un número limitado de turnos y reduzca de forma notable la frecuencia de alineación manual.

\paragraph{}
El marco DSPy utiliza el latent cost \(C\) como puente para enlazar el search trace, la salida de APAC y el solver, reduciendo la dimensionalidad del problema de combinatorial constraint hasta obtener una trayectoria de optimización diferenciable. El Nano optics benchmark demostró que este pipeline puede atender al mismo tiempo el desempeño y la validez del proceso de fabricación.

\paragraph{}
Lo más destacado de esta clase es que enlaza search, agent y solver en un solo flujo de información, en lugar de tratarlos como módulos independientes: cada salida incluye un time-stamped trace, un judge reward y un solver plan, y en conjunto forman un audit log completo.

\subsection{Resumen de la sección}
Al repasar esta sección se observa que, desde el search trace hasta el multi-turn agent y de ahí al DSPy latent solver, toda la arquitectura gira en torno a la interpretabilidad, la supervisión estructurada y la observabilidad de la ingeniería, formando un bucle cerrado reproducible y depurable.

